% --- Archon physics formalization source begin ---
% archon:physics
% archon:covers PhyXMiniProblems/problem_phyx_mini_0130.lean
% archon:source-report reports/phyx_mini/problem_phyx_mini_0130.source.json

\chapter{Physics problem phyx\_mini\_0130}
\label{ch:PhyXMiniProblems_problem_phyx_mini_0130}

\paragraph{Problem source.}
\#\# Physical scenario

A nearsighted eye has near and far points of \textbackslash{}(12\textbackslash{},\textbackslash{}mathrm\{cm\}\textbackslash{}) and \textbackslash{}(17\textbackslash{},\textbackslash{}mathrm\{cm\}\textbackslash{}), respectively.Assume that the lens is \textbackslash{}(2.0\textbackslash{},\textbackslash{}mathrm\{cm\}\textbackslash{}) from the eye (typical for eyeglasses).

\#\# Question

What lens power is needed for this person to see distant objects clearly?

\#\# Answer choices

- A: A: -5.3D

- B: B: -6.3D

- C: C: -6.7D

- D: D: -5.6D

\#\# Figure caption (auxiliary; use the image as primary evidence)

The image is a diagram illustrating how light is focused on the eye using a lens. 

- **Objects and Scenes**: 
  - There is an eye on the right with focus lines directed toward it.
  - A convex lens is placed in front of the eye.
  - Two points labeled "O" and "I" are on the left side of the lens, representing the object and image, respectively.

- **Relationships**:
  - The object point "O" emits light rays, which converge at point "I" after passing through the lens, directed towards the eye.
  - The image point "I" is closer to the lens than the object point "O".

- **Text and Measurements**:
  - The distance between point "I" and the lens is marked as "12 cm" and is labeled as the "Near point".

\#\# Dataset metadata

Geometrical Optics; Physical Model Grounding Reasoning, Multi-Formula Reasoning

\paragraph{Recorded answer/context.}
C: C: -6.7D

\paragraph{Figure/image path.}
/root/proposal\_for\_physic/science-mango/phyx\_mini\_run/phyx\_data/test\_image/130.png

\paragraph{Formalization target.}
create a compiling Lean file with sorry bodies at `PhyXMiniProblems/problem\_phyx\_mini\_0130.lean`. The Lean declarations must preserve the physical quantities, dimensions or dimensional roles, figure labels, governing-law hypotheses, and final relation expressed by this problem.
Use Mathlib/Physlib names found through LeanExplore where available. If a physics API is missing, introduce faithful local abstractions rather than scalar placeholder aliases.

\begin{theorem}[Physics formalization target]
\label{thm:physics:phyx_mini_0130:target}
\lean{PhyXMiniProblems.ProblemPhyXMini0130.problem_phyx_mini_0130}
\uses{def:physics:phyx-mini-0130:phyxminiproblems-problemphyxmini0130-powerindiopters, def:physics:phyx-mini-0130:phyxminiproblems-problemphyxmini0130-nearsightedcorrectionsetup, def:physics:phyx-mini-0130:phyxminiproblems-problemphyxmini0130-hasstateddistancereadouts, def:physics:phyx-mini-0130:phyxminiproblems-problemphyxmini0130-hasphysicaldistanceordering, def:physics:phyx-mini-0130:phyxminiproblems-problemphyxmini0130-matchessourcefigure, def:physics:phyx-mini-0130:phyxminiproblems-problemphyxmini0130-viewsdistantobject, def:physics:phyx-mini-0130:phyxminiproblems-problemphyxmini0130-formsvirtualimageatunaidedfarpoint, def:physics:phyx-mini-0130:phyxminiproblems-problemphyxmini0130-obeysthinlensequation, def:physics:phyx-mini-0130:phyxminiproblems-problemphyxmini0130-obeysopticalpowerlaw, def:physics:phyx-mini-0130:phyxminiproblems-problemphyxmini0130-lenskindagreeswithfocalsign, def:physics:phyx-mini-0130:phyxminiproblems-problemphyxmini0130-answerpowerdiopters, def:physics:phyx-mini-0130:phyxminiproblems-problemphyxmini0130-isnearesttenthreadout}
The assigned autoformalize agent should translate this physics problem into Lean declarations in the covered file, with theorem and lemma proof bodies written as `by sorry`.
\end{theorem}
\begin{proof}
This is an autoformalization task, not a proof task. Produce statements that can later be proved by the physics prover without weakening the physical model.
\end{proof}

% --- Archon declaration topology begin ---
\section{Declaration topology}
\begin{definition}[Length Quantity]
  \label{def:physics:phyx-mini-0130:phyxminiproblems-problemphyxmini0130-lengthquantity}
  \lean{PhyXMiniProblems.ProblemPhyXMini0130.LengthQuantity}
  A signed physical length with coherent readouts in every unit system
\end{definition}
\begin{proof}
  This is the declared model definition of Length Quantity.
\end{proof}
\begin{definition}[Optical Power Quantity]
  \label{def:physics:phyx-mini-0130:phyxminiproblems-problemphyxmini0130-opticalpowerquantity}
  \lean{PhyXMiniProblems.ProblemPhyXMini0130.OpticalPowerQuantity}
  A physical lens power, whose SI readout is measured in diopters
\end{definition}
\begin{proof}
  This is the declared model definition of Optical Power Quantity.
\end{proof}
\begin{definition}[length In Meters]
  \label{def:physics:phyx-mini-0130:phyxminiproblems-problemphyxmini0130-lengthinmeters}
  \lean{PhyXMiniProblems.ProblemPhyXMini0130.lengthInMeters}
  \uses{def:physics:phyx-mini-0130:phyxminiproblems-problemphyxmini0130-lengthquantity}
  The SI scalar readout of a physical length, in metres
\end{definition}
\begin{proof}
  This is the declared model definition of length In Meters.
\end{proof}
\begin{definition}[power In Diopters]
  \label{def:physics:phyx-mini-0130:phyxminiproblems-problemphyxmini0130-powerindiopters}
  \lean{PhyXMiniProblems.ProblemPhyXMini0130.powerInDiopters}
  \uses{def:physics:phyx-mini-0130:phyxminiproblems-problemphyxmini0130-opticalpowerquantity}
  The SI scalar readout of optical power, in inverse metres (diopters)
\end{definition}
\begin{proof}
  This is the declared model definition of power In Diopters.
\end{proof}
\begin{definition}[Thin Lens Kind]
  \label{def:physics:phyx-mini-0130:phyxminiproblems-problemphyxmini0130-thinlenskind}
  \lean{PhyXMiniProblems.ProblemPhyXMini0130.ThinLensKind}
  Optical type of a paraxial thin lens
\end{definition}
\begin{proof}
  This is the declared model definition of Thin Lens Kind.
\end{proof}
\begin{definition}[Image Kind]
  \label{def:physics:phyx-mini-0130:phyxminiproblems-problemphyxmini0130-imagekind}
  \lean{PhyXMiniProblems.ProblemPhyXMini0130.ImageKind}
  Whether rays themselves meet at an image or only appear to come from it
\end{definition}
\begin{proof}
  This is the declared model definition of Image Kind.
\end{proof}
\begin{definition}[Figure Point]
  \label{def:physics:phyx-mini-0130:phyxminiproblems-problemphyxmini0130-figurepoint}
  \lean{PhyXMiniProblems.ProblemPhyXMini0130.FigurePoint}
  The four labels whose left-to-right order is visible in the source figure
\end{definition}
\begin{proof}
  This is the declared model definition of Figure Point.
\end{proof}
\begin{definition}[Object Location]
  \label{def:physics:phyx-mini-0130:phyxminiproblems-problemphyxmini0130-objectlocation}
  \lean{PhyXMiniProblems.ProblemPhyXMini0130.ObjectLocation}
  \uses{def:physics:phyx-mini-0130:phyxminiproblems-problemphyxmini0130-lengthquantity}
  An object's axial location for the paraxial lens equation. Optical infinity is kept as a genuine alternative to a finite physical distance
\end{definition}
\begin{proof}
  This is the declared model definition of Object Location.
\end{proof}
\begin{definition}[object Vergence In Per Meters]
  \label{def:physics:phyx-mini-0130:phyxminiproblems-problemphyxmini0130-objectvergenceinpermeters}
  \lean{PhyXMiniProblems.ProblemPhyXMini0130.objectVergenceInPerMeters}
  \uses{def:physics:phyx-mini-0130:phyxminiproblems-problemphyxmini0130-lengthinmeters, def:physics:phyx-mini-0130:phyxminiproblems-problemphyxmini0130-objectlocation}
  Object vergence in inverse metres; parallel rays from infinity have zero vergence
\end{definition}
\begin{proof}
  This is the declared model definition of object Vergence In Per Meters.
\end{proof}
\begin{definition}[Clear Vision Range]
  \label{def:physics:phyx-mini-0130:phyxminiproblems-problemphyxmini0130-clearvisionrange}
  \lean{PhyXMiniProblems.ProblemPhyXMini0130.ClearVisionRange}
  \uses{def:physics:phyx-mini-0130:phyxminiproblems-problemphyxmini0130-lengthquantity}
  The interval of object distances that the unaided eye can focus clearly
\end{definition}
\begin{proof}
  This is the declared model definition of Clear Vision Range.
\end{proof}
\begin{definition}[Thin Corrective Lens]
  \label{def:physics:phyx-mini-0130:phyxminiproblems-problemphyxmini0130-thincorrectivelens}
  \lean{PhyXMiniProblems.ProblemPhyXMini0130.ThinCorrectiveLens}
  \uses{def:physics:phyx-mini-0130:phyxminiproblems-problemphyxmini0130-lengthquantity, def:physics:phyx-mini-0130:phyxminiproblems-problemphyxmini0130-opticalpowerquantity, def:physics:phyx-mini-0130:phyxminiproblems-problemphyxmini0130-thinlenskind}
  A thin corrective lens with physically distinct focal length and power
\end{definition}
\begin{proof}
  This is the declared model definition of Thin Corrective Lens.
\end{proof}
\begin{definition}[Nearsighted Correction Setup]
  \label{def:physics:phyx-mini-0130:phyxminiproblems-problemphyxmini0130-nearsightedcorrectionsetup}
  \lean{PhyXMiniProblems.ProblemPhyXMini0130.NearsightedCorrectionSetup}
  \uses{def:physics:phyx-mini-0130:phyxminiproblems-problemphyxmini0130-lengthquantity, def:physics:phyx-mini-0130:phyxminiproblems-problemphyxmini0130-imagekind, def:physics:phyx-mini-0130:phyxminiproblems-problemphyxmini0130-figurepoint, def:physics:phyx-mini-0130:phyxminiproblems-problemphyxmini0130-objectlocation, def:physics:phyx-mini-0130:phyxminiproblems-problemphyxmini0130-clearvisionrange, def:physics:phyx-mini-0130:phyxminiproblems-problemphyxmini0130-thincorrectivelens}
  The quantities and labels needed to model both the given eye and its spectacle correction. The signed distant-image distance is negative when the image lies on the incoming-light side of the lens
\end{definition}
\begin{proof}
  This is the declared model definition of Nearsighted Correction Setup.
\end{proof}
\begin{definition}[Has Stated Distance Readouts]
  \label{def:physics:phyx-mini-0130:phyxminiproblems-problemphyxmini0130-hasstateddistancereadouts}
  \lean{PhyXMiniProblems.ProblemPhyXMini0130.HasStatedDistanceReadouts}
  \uses{def:physics:phyx-mini-0130:phyxminiproblems-problemphyxmini0130-lengthinmeters, def:physics:phyx-mini-0130:phyxminiproblems-problemphyxmini0130-nearsightedcorrectionsetup}
  The near point, far point, and spectacle-lens offset stated in the problem
\end{definition}
\begin{proof}
  This is the declared model definition of Has Stated Distance Readouts.
\end{proof}
\begin{definition}[Has Physical Distance Ordering]
  \label{def:physics:phyx-mini-0130:phyxminiproblems-problemphyxmini0130-hasphysicaldistanceordering}
  \lean{PhyXMiniProblems.ProblemPhyXMini0130.HasPhysicalDistanceOrdering}
  \uses{def:physics:phyx-mini-0130:phyxminiproblems-problemphyxmini0130-lengthinmeters, def:physics:phyx-mini-0130:phyxminiproblems-problemphyxmini0130-nearsightedcorrectionsetup}
  The given eye distances are positive and ordered as a nearsighted range
\end{definition}
\begin{proof}
  This is the declared model definition of Has Physical Distance Ordering.
\end{proof}
\begin{definition}[Matches Source Figure]
  \label{def:physics:phyx-mini-0130:phyxminiproblems-problemphyxmini0130-matchessourcefigure}
  \lean{PhyXMiniProblems.ProblemPhyXMini0130.MatchesSourceFigure}
  \uses{def:physics:phyx-mini-0130:phyxminiproblems-problemphyxmini0130-lengthinmeters, def:physics:phyx-mini-0130:phyxminiproblems-problemphyxmini0130-figurepoint, def:physics:phyx-mini-0130:phyxminiproblems-problemphyxmini0130-nearsightedcorrectionsetup}
  Primary-image readout: O, I, the concave spectacle lens, and the eye occur from left to right. The dashed backward ray extension identifies I as a virtual image, and the marked I--eye separation is the 12 cm near point
\end{definition}
\begin{proof}
  This is the declared model definition of Matches Source Figure.
\end{proof}
\begin{definition}[Views Distant Object]
  \label{def:physics:phyx-mini-0130:phyxminiproblems-problemphyxmini0130-viewsdistantobject}
  \lean{PhyXMiniProblems.ProblemPhyXMini0130.ViewsDistantObject}
  \uses{def:physics:phyx-mini-0130:phyxminiproblems-problemphyxmini0130-nearsightedcorrectionsetup}
  The object whose distant-vision correction is requested is at infinity
\end{definition}
\begin{proof}
  This is the declared model definition of Views Distant Object.
\end{proof}
\begin{definition}[Forms Virtual Image At Unaided Far Point]
  \label{def:physics:phyx-mini-0130:phyxminiproblems-problemphyxmini0130-formsvirtualimageatunaidedfarpoint}
  \lean{PhyXMiniProblems.ProblemPhyXMini0130.FormsVirtualImageAtUnaidedFarPoint}
  \uses{def:physics:phyx-mini-0130:phyxminiproblems-problemphyxmini0130-lengthinmeters, def:physics:phyx-mini-0130:phyxminiproblems-problemphyxmini0130-nearsightedcorrectionsetup}
  Prescription criterion for myopia: the corrective lens makes parallel rays appear to originate at the unaided far point. Subtracting the lens-to-eye offset converts the far-point distance from an eye-based to a lens-based distance, and the virtual-image sign is negative
\end{definition}
\begin{proof}
  This is the declared model definition of Forms Virtual Image At Unaided Far Point.
\end{proof}
\begin{definition}[Obeys Thin Lens Equation]
  \label{def:physics:phyx-mini-0130:phyxminiproblems-problemphyxmini0130-obeysthinlensequation}
  \lean{PhyXMiniProblems.ProblemPhyXMini0130.ObeysThinLensEquation}
  \uses{def:physics:phyx-mini-0130:phyxminiproblems-problemphyxmini0130-lengthinmeters, def:physics:phyx-mini-0130:phyxminiproblems-problemphyxmini0130-objectvergenceinpermeters, def:physics:phyx-mini-0130:phyxminiproblems-problemphyxmini0130-nearsightedcorrectionsetup}
  The signed paraxial thin-lens equation 1/f = Vₒ + 1/dᵢ
\end{definition}
\begin{proof}
  This is the declared model definition of Obeys Thin Lens Equation.
\end{proof}
\begin{definition}[Obeys Optical Power Law]
  \label{def:physics:phyx-mini-0130:phyxminiproblems-problemphyxmini0130-obeysopticalpowerlaw}
  \lean{PhyXMiniProblems.ProblemPhyXMini0130.ObeysOpticalPowerLaw}
  \uses{def:physics:phyx-mini-0130:phyxminiproblems-problemphyxmini0130-lengthinmeters, def:physics:phyx-mini-0130:phyxminiproblems-problemphyxmini0130-powerindiopters, def:physics:phyx-mini-0130:phyxminiproblems-problemphyxmini0130-thincorrectivelens}
  Governing relation P = 1/f between optical power and focal length
\end{definition}
\begin{proof}
  This is the declared model definition of Obeys Optical Power Law.
\end{proof}
\begin{definition}[Lens Kind Agrees With Focal Sign]
  \label{def:physics:phyx-mini-0130:phyxminiproblems-problemphyxmini0130-lenskindagreeswithfocalsign}
  \lean{PhyXMiniProblems.ProblemPhyXMini0130.LensKindAgreesWithFocalSign}
  \uses{def:physics:phyx-mini-0130:phyxminiproblems-problemphyxmini0130-lengthinmeters, def:physics:phyx-mini-0130:phyxminiproblems-problemphyxmini0130-thincorrectivelens}
  The focal-length sign agrees with the physical type of the thin lens
\end{definition}
\begin{proof}
  This is the declared model definition of Lens Kind Agrees With Focal Sign.
\end{proof}
\begin{definition}[Answer Choice]
  \label{def:physics:phyx-mini-0130:phyxminiproblems-problemphyxmini0130-answerchoice}
  \lean{PhyXMiniProblems.ProblemPhyXMini0130.AnswerChoice}
  Labels of the four answer choices printed with the problem
\end{definition}
\begin{proof}
  This is the declared model definition of Answer Choice.
\end{proof}
\begin{definition}[answer Power Diopters]
  \label{def:physics:phyx-mini-0130:phyxminiproblems-problemphyxmini0130-answerpowerdiopters}
  \lean{PhyXMiniProblems.ProblemPhyXMini0130.answerPowerDiopters}
  \uses{def:physics:phyx-mini-0130:phyxminiproblems-problemphyxmini0130-answerchoice}
  The displayed optical-power readout of each answer, in diopters
\end{definition}
\begin{proof}
  This is the declared model definition of answer Power Diopters.
\end{proof}
\begin{definition}[Is Nearest Tenth Readout]
  \label{def:physics:phyx-mini-0130:phyxminiproblems-problemphyxmini0130-isnearesttenthreadout}
  \lean{PhyXMiniProblems.ProblemPhyXMini0130.IsNearestTenthReadout}
  A reported value is an integer number of tenths within half a tenth of the exact value
\end{definition}
\begin{proof}
  This is the declared model definition of Is Nearest Tenth Readout.
\end{proof}
\begin{lemma}[signed Distant Image Distance eq neg three twentieth]
  \label{lem:physics:phyx-mini-0130:phyxminiproblems-problemphyxmini0130-signeddistantimagedistance-eq-neg-three-twentieth}
  \lean{PhyXMiniProblems.ProblemPhyXMini0130.signedDistantImageDistance_eq_neg_three_twentieth}
  \uses{def:physics:phyx-mini-0130:phyxminiproblems-problemphyxmini0130-lengthinmeters, def:physics:phyx-mini-0130:phyxminiproblems-problemphyxmini0130-nearsightedcorrectionsetup, def:physics:phyx-mini-0130:phyxminiproblems-problemphyxmini0130-hasstateddistancereadouts, def:physics:phyx-mini-0130:phyxminiproblems-problemphyxmini0130-formsvirtualimageatunaidedfarpoint}
  The required virtual image is 15 cm to the left of the spectacle lens
\end{lemma}
\begin{proof}
  The conclusion follows from the governing relations and figure readouts named in the statement of signed Distant Image Distance eq neg three twentieth.
\end{proof}
\begin{lemma}[reciprocal Focal Length eq neg twenty thirds]
  \label{lem:physics:phyx-mini-0130:phyxminiproblems-problemphyxmini0130-reciprocalfocallength-eq-neg-twenty-thirds}
  \lean{PhyXMiniProblems.ProblemPhyXMini0130.reciprocalFocalLength_eq_neg_twenty_thirds}
  \uses{def:physics:phyx-mini-0130:phyxminiproblems-problemphyxmini0130-lengthinmeters, def:physics:phyx-mini-0130:phyxminiproblems-problemphyxmini0130-nearsightedcorrectionsetup, def:physics:phyx-mini-0130:phyxminiproblems-problemphyxmini0130-hasstateddistancereadouts, def:physics:phyx-mini-0130:phyxminiproblems-problemphyxmini0130-viewsdistantobject, def:physics:phyx-mini-0130:phyxminiproblems-problemphyxmini0130-formsvirtualimageatunaidedfarpoint, def:physics:phyx-mini-0130:phyxminiproblems-problemphyxmini0130-obeysthinlensequation}
  The distant-object lens law fixes the reciprocal focal length at -20/3 m⁻¹
\end{lemma}
\begin{proof}
  The conclusion follows from the governing relations and figure readouts named in the statement of reciprocal Focal Length eq neg twenty thirds.
\end{proof}
% --- Archon declaration topology end ---
% --- Archon physics formalization source end ---
